% chapters/ch15.tex —— 第 15 章
% 体例：具体失败 → 设计问题 → 机制解释 → 对照实验 → 条件变化 → 选择原则
% 数字身份纪律：
%   - E01 / E02 / E03 / E03c 的运行读数一律 \measured{}；
%   - 外部论文报告值一律 \cited{}；
%   - 由上述数字算出的倍数、差值与比例用 \calcnum{}。
% E03c 第 10 臂 fit_domain_cross 已于 2026-09-28 用 --resume-latest 续跑补齐（原阻塞于换模型后端余额 HTTP 402），
% 本章不再写"待补"；跨模型读数已出但仍区间跨 0，P3b关注拟合目标与读取者条件。
\chapter{试验审计与迁移设计}
\label{ch15}

\chaptergoal{沿既有工件重走试验审计，区分可复核的观察、尚未隔离的组件作用与下一轮迁移设计。}
\chapterprops{P3b（拟合目标怎样影响跨读取者迁移）。
同读取者的方案差为\measured{$0.0000$}、区间跨零；替代读取者的方案差为\measured{$+0.0625$}、
区间\measured{$[-0.0625, +0.1750]$}同样跨零。该观察用于提出拟合目标与读取者的交互问题。}
\chapterfig{图~\ref{fig:ch15-reader}：两种写法 × 两个读取者的四格}
\chapterpitfall{两个无记忆静态基线换读取者后，通过率从 \measured{0.6875} 掉到 \measured{0.2625}——
有人把它写成“经验失效”，有人写成“模型更弱”，两个数都是真的}

% ---------- 1. 具体失败 ----------
\begin{failurebox}
E03c 的拆分 B 里有一行最容易读错的数字：同一个持久物 \texttt{fit\_solver}，
在同模型读取者手上 回合通过率 是 \measured{0.7000}，换一个读取者之后是 \measured{0.3625}。
只看这两个数，容易把总表现下降误读成经验增量失效。

但同一张表里还有一行必须一起看的：\texttt{base\_same} \measured{0.6875} 对
\texttt{base\_cross} \measured{0.2625}——\textbf{不新增跨任务记忆的静态基线，
换一次读取者也掉了 \calcnum{42.5} 个百分点}。
换模型后，经验臂减去对应静态基线，\texttt{fit\_solver} 的局部收益是
\measured{+0.1000}（区间 \measured{$[-0.0125, +0.2125]$}），点估计为正，区间仍跨零。

同一个 \measured{0.3625}，减去源读取者基线得到负差，减去目标读取者基线则得到正差；两次相减回答不同问题。
这说明读数必须结合相应基线：\textbf{有没有同时报告换读取者之后的静态基线}。
\end{failurebox}

\designquestion{同一份归档材料究竟能回答哪个问题？本章沿配置、经验、逐次轨迹、汇总和解释走查E01、E03与E05v2，练习发现基线错误、记录冲突和提前失败。组件叠加、跨仓库、跨环境及真实恢复是后续设计，尚无本书实测。}

读者可依次完成四步离线审计：先按附录D核对运行身份与哈希，再按第5章核对E03逐次判定，随后运行本章案例索引复核E05v2的六个序列，最后写出缺少哪项干预才能回答组件归因问题。这里交付的是可跟随的审计流程；完整集成实验仍需新协议和另行授权。

% ---------- 2. 静态基线与错误谱 ----------
\section{从已有基线与错误谱开始}
\label{sec:ch15-baseline}

\begin{sloppypar}
任何“收益”的说法都要先有一个不动的参照物。本书的参照物不是“没有 Agent”，
而是\textbf{一个保留任务工具、但不发布跨任务经验的系统}——它有工具、能重跑、能看反馈，
只是不携带本试验中新形成的跨任务经验。
\end{sloppypar}

\begin{sloppypar}
在这条基线上，E01 的四格给出了全套读数：同一批封存题、同样的 episode 配额、
相同请求模型与回合上限，但起始经验来源、旁路生成及实际用量仍需分别核对，不能称为只改变一个组件。
\end{sloppypar}

\begin{table}[htbp]
\centering
\caption{E01四种方案：回合通过率、通过次数与主回合计时和}
\label{tab:ch15-baseline}
\begin{small}
\begin{tabular}{>{\raggedright\arraybackslash}p{1.2cm}>{\raggedright\arraybackslash}p{4.0cm}ccc}
\toprule
格 & 改动留在哪一层 & 回合通过率 & 通过 / 总运行 & 回合计时和（秒） \\
\midrule
格 1 & 只留在回合内，重复取均值 & \measured{0.7375} & \measured{59}/\measured{80} & \measured{369.6} \\
格 2 & 持久物外部构建，全程只读 & \measured{0.7250} & \measured{58}/\measured{80} & \measured{417.7} \\
格 3 & 每道题单独适配，用完即弃 & \measured{0.7625} & \measured{61}/\measured{80} & \measured{411.3} \\
格 4 & 跨任务持久更新 & \measured{0.7875} & \measured{63}/\measured{80} & \measured{407.1} \\
\bottomrule
\end{tabular}
\end{small}
\end{table}

{\footnotesize 通过次数的分母只含\measured{80}个封存评价回合；计时列汇总每臂\measured{88}个主回合，包含\measured{8}个训练回合，不含旁路生成等待。两列对应不同范围，不能将计时列当成80个评价回合的总耗时。}


\begin{sloppypar}
错误谱可以同时按题和按族检查。八个族每族 \measured{2} 道配对题（格 4 减格 1），
点估计从 \measured{$-0.3000$}（\texttt{recursionbase}）到 \measured{+0.4000}（\texttt{mutabledefault}）都有：
\texttt{classstate} \measured{$-0.1000$}、\texttt{dictmiss} \measured{+0.1000}、
\texttt{errorcontract} \measured{$0.0000$}、\texttt{sortkey} \measured{+0.3000}、
\texttt{offbyone} \measured{$-0.1000$}、\texttt{strindex} \measured{+0.1000}。
每族两道题只能提供这个有限样本的描述，不能可靠代表整个族；这张谱用于提出后续问题，不能直接当作已知的总体方差。
它真正的用处是提醒一件事：总体为正完全可能由一两个族撑起来。
\end{sloppypar}

\begin{pitfallbox}
把“静态基线”设成“没有记忆的同一个 Agent”而不是“没有 Agent”，是本书一直坚持的做法。
前者为评价新增经验提供参照；是否识别写入本身的作用，还取决于其他条件是否匹配。后者则改变了是否使用Agent这一条件。
两种参照回答不同问题；本书没有用对应对照测出二者固定相差多少，不能填入一个通用百分点差。
\end{pitfallbox}

% ---------- 3. 逐步叠加 ----------
\subsection{正负点估计与恢复证据缺口}
\label{sec:ch15-cases}

\begin{sloppypar}
下面区分观测到的点估计、尚未判定的效应与缺失的恢复证据。分类用于阅读已有记录，不把点值的正负自动转成成功或有害的科学判决。
\end{sloppypar}

\begin{itemize}
  \item \begin{sloppypar}\textbf{正向点估计}：E03拆分B中，同模型两种拟合方案均为\measured{0.7000}，相对静态基线\measured{0.6875}均为\measured{+0.0125}；对应区间\measured{$[-0.1125,+0.1250]$}与\measured{$[-0.1500,+0.1625]$}跨零。它们是探索性读数，不能仅因符号为正就记为已经证明有效的更新。\end{sloppypar}
  \item \textbf{负向点估计}：E02五个有记忆臂的成绩均低于无记忆基线\measured{0.8125}，自评臂为\measured{0.6125}，差\measured{$-0.2000$}。各臂不确定性与操作定义不同，不能把五个点值统一当作“记忆无效”的结论。
  \item \textbf{矛盾反馈案例}：噪声臂在\measured{8}条反馈中翻转\measured{2}条判定标签，但原执行输出仍保留。相对基线点估计\measured{$-0.1500$}，区间\measured{$[-0.3125,+0.0125]$}；题级均值有\measured{8}题下降、\measured{3}题上升。这没有隔离纯标签噪声，也不保证最终发布了错误经验；噪声臂与测试判定臂还分别训练并生成不同经验，这一方案比较未隔离两条标签翻转的净效应；未部署报警器的记录也不能估计报警漏检率。
  \item \textbf{恢复证据缺口}：本书尚无E12真实模型恢复率。相关综述建议报告恢复成功率与成本\cite{coding-survey-2608}，并不意味着所有已有方法都没有数值。新版E12的文件恢复离线检查已经通过，仍不替代真实模型复评。
\end{itemize}

\begin{sloppypar}
E03c求解器拟合方案有\measured{1}/\measured{16}题出现相对各自静态基线的点估计由正转负，问题域拟合为\measured{2}/\measured{16}。这些现象提示适用条件值得保存，不能从少量重复直接判定每道题的真实效应反号。第5章的双字段设计是一种表达范围的方式，实际能否改善迁移还需要自己的对照，不能由案例直接推出。
\end{sloppypar}

% ---------- 8. 换读取者 ----------
\section{组件归因需要怎样的共同起点}
\label{sec:ch15-stack}

\begin{sloppypar}
第 2 章把系统写成 $A_k=(\theta,P,M,L,G;\mathcal V)$。
完整集成比较可以逐步开放这五类可变组件，检查每一步的资源与质量变化；本书尚未运行这条完整版本链。表中的记忆操作是记忆组件的控制方式，单列说明，并非第六类持久对象。
下表按本书已有的实测状态逐行登记——其中的空项标明证据缺口，不表示相应机制无效。
\end{sloppypar}

\begin{table}[htbp]
\centering
\caption{五类持久组件的叠加顺序与本书的实测状态}
\label{tab:ch15-stack}
\begin{footnotesize}
\begin{tabular}{>{\raggedright\arraybackslash}p{1.6cm}>{\raggedright\arraybackslash}p{1.0cm}>{\raggedright\arraybackslash}p{4.4cm}>{\raggedright\arraybackslash}p{3.0cm}>{\raggedright\arraybackslash}p{2.2cm}}
\toprule
组件 & 章 & 打开它意味着什么 & 本书的落点 & 实测状态 \\
\midrule
$M$ 经验记忆 & 5 & 经验被写下来并被下次检索到 & E01 格 4、E03 拆分 A/B & 有实测；A基线暂停 \\
\midrule
记忆操作 & 6 & 写 / 留 / 取 / 删本身可学习 & 准入、组织、触发与操作策略（第 6 章） & 未跑对照 \\
\midrule
$P$ 指令与策略 & 7 & 提示与持久策略被改写 & 第 7 章机制章 & 未跑对照 \\
\midrule
$L$ 工具与技能 & 8 & 新技能被固化进库并复用 & 第 8 章机制章 & E05 仅工程试跑，无判决 \\
\midrule
$G$ 工作流与结构 & 9 & 角色与流程本身被重排 & 第 9 章机制章 & 未跑对照 \\
\midrule
$\theta$ 参数 & 10、11 & 经验进入权重 & 第 10、11 章 & 未跑对照 \\
\bottomrule
\end{tabular}
\end{footnotesize}
\end{table}

\begin{sloppypar}
记忆相关实验已有多臂记录，技能已有E05 v2工程小试；其余组件缺相应实测。已有运行尚未完成组件归因，本章因此提供机制走查和下一步设计，不声称整套系统已验证。
\end{sloppypar}

% ---------- 4. 等预算比较 ----------
\subsection{审计回合配额与组件归因}
\label{sec:ch15-ablation}

\begin{sloppypar}
四格的评价回合配额相同，均为\measured{16}题乘\measured{5}次；相同上限不等于相同实际总资源。历史主回合计时和从\measured{369.6}秒到\measured{417.7}秒，较大值高约\calcnum{13.0\%}，但这不是包含全部准备与旁路的端到端墙钟，也不是美元费用排序。完整成本范围和重试策略的比较见第~\ref{ch03}章与附录~\ref{app:cost}。\end{sloppypar}

\begin{sloppypar}
在这个预算下，格 4 相对格 1 的点估计是 \measured{+0.0500}，
区间 \measured{$[-0.1000, +0.1750]$}；
相对格 2 是 \measured{+0.0625}，区间 \measured{$[-0.1375, +0.2750]$}。
\textbf{两个区间都跨 0}，这既不证明差异为零，也不解除完整预算与来源差异的设计限制；当前保留探索性读数。
\end{sloppypar}

\begin{sloppypar}
更有信息量的是格 2 与格 4 的对照，它在读写设置上接近 P1b 的问题：
格 2 是“有持久物但全程只读”（\measured{0.7250}），格 4 是“有持久物且允许写”（\measured{0.7875}）。
两者的差\measured{+0.0625}描述两个完整方案；起始经验来源不同，不能只归因于是否允许写入。
它提醒我们分别检查初始材料和写入规则；格2与格4的现有差异尚未隔离持续写入的增量。
\textbf{格 1 与格 4 的差里混着“有东西可读”和“允许写”两件事}。
\end{sloppypar}

\begin{sloppypar}
本书还保存一个用于实现核对的字段：\texttt{persistent\_writes} 字段。
静态基线的轨迹里这个字段一旦非空，运行器直接抛异常。
这是对被记录的持久写入字段所作的运行后断言，不是阻止任意文件或外部副作用的文件系统安全机制；缺失的写入记录也不能由空字段排除。
\end{sloppypar}

% ---------- 5. 跨族与跨仓库 ----------
\subsection{把组件连起来，先固定一条版本链}
\label{sec:ch15-versionchain}

假设下一阶段要在同一开发助手中同时加入经验记忆、策略提示和代码参考。最直接的办法是一次全部打开，看总成功率是否提高；这能回答整套配置是否值得使用，却难以说明各组件的作用。为了让集成案例具有解释力，应先保存共同起点，并为每一次变化生成一个可识别版本。

初始版本可以包含固定模型、工具、回合预算和任务接口，但不跨任务写入经验。第一个候选只允许把来源记录存入记忆，第二个在同一记忆条件上允许调整策略，第三个再加入代码参考。每个版本都要绑定持久物快照、提示和工具定义。这里是未来组合试验的设计，本书现有 E01、E03、E05 来自不同运行，不能直接代替这条受控版本链。

逐步加入只得到条件性增量。记忆在没有技能时的作用，与已有技能后的作用可能不同；策略也可能改变记忆检索次数。如果只按一种顺序增加组件，最后加入的组件得到的是此前配置之上的边际效果，而非独立于其他组件的固定贡献。逐步叠加仍可回答工程问题，其解释应限定为所选顺序和共同起点下的增量。

可以在集成后再做移除比较：在同一个完整版本上关闭一项机制，观察任务、成本和失效类型的变化。但关闭方式要保持接口可运行。例如删除技能模块后仍让提示要求调用它，会制造一个接口损坏条件；这种下降不能代表技能本身的正常贡献。消融应定义一个有效替代路径，而不是故意让系统缺少必需入口。

组件之间还可能存在协同。以教学记号 $Y(M,L)$ 表示记忆和技能开关下的表现，交互差可以写成
\begin{equation}
 I=Y(1,1)-Y(1,0)-Y(0,1)+Y(0,0).
 \label{eq:ch15-interaction}
\end{equation}
它需要四个实际条件，不能只用基线与完整系统两格估计。若 $I$ 为正，表示在所定义尺度上组合增量超过两个单独增量之和；这仍是特定指标、任务和预算下的关系，不是组件普遍互补的定律。

用人为输入作一次读数练习：基线\calcnum{0.50}，单独记忆\calcnum{0.55}，单独技能\calcnum{0.60}，组合\calcnum{0.62}。组合表现最高，但交互差为\calcnum{$-0.03$}。这不是矛盾：组合仍有收益，只是两个组件的增量部分重叠。若只展示最高的一格，就会错过决定下一步投入的重要信息。

\subsection{共同起点需要包括环境，不只是模型名}
\label{sec:ch15-commonstart}

模型请求名相同，只固定了比较条件中的一部分。代码仓库、依赖、缓存、工作目录、工具返回格式、当前时间和初始持久物都可能影响行为。若一条路线继承了前一条路线留下的文件，另一条从干净状态开始，比较就已经改变。集成试验应从可核对的环境快照启动每个需要独立的条件。

本书短程序任务可以通过任务源文件与测试文件重建局部环境；更复杂仓库则还涉及依赖安装、测试数据和外部服务。任务准备失败不应被默默转换成求解失败，也不能从分母里无记录删除。可以预先定义环境不可用的处理规则，并将基础设施结果与算法任务指标分别保存。

执行顺序也可能产生影响。同一服务在不同时间的延迟或返回行为可能变化；多个条件共用缓存，会改变费用和速度。随机化或交错运行可以减轻某些时序混淆，但仍需与资源和服务约束匹配。如果只能固定顺序，应保留这一限制，不把顺序差当成组件效应。

共同起点不要求所有后续状态永远相同。更新机制本来就会让持久物和策略分化；需要固定的是产生分化之前的条件，以及允许分化的规则。一个方法成功更多，因此积累了更多经过验收的经验，可以是闭环系统的真实效果；若研究只想比较同一经验的表示方式，就应改用冻结经验包。两类实验的起点要求相似，信息流要求不同。

可以在每次运行开始时输出一份简洁指纹：任务清单、输入文件哈希、配置、模型请求名与可得响应身份、持久物标识、工具版本和时间范围。指纹中缺少服务端固定版本时应写明，而不能用客户端包版本冒充模型身份。第3章和附录D讨论的记录纪律在集成阶段更重要，因为影响结果的组件更多了。

\subsection{组合系统的发布记录应该长什么样}
\label{sec:ch15-release}

一次集成候选准备发布时，最有用的摘要不是“所有测试绿色”，而是说明它改变了什么、在哪些任务上怎样测过、目前覆盖到哪里。可以从共同起点版本出发，列出新增持久物、策略与工具，随后报告相应任务结果、资源范围、失败类型和恢复条件。细节留在可定位工件中，正文保留决定接受与否的依据。

若候选只在开发范围表现更好，可以继续受限测试，却不应把开发结果当作独立迁移证明。若最终范围出现负向点值但精度不足，可以报告不确定并根据业务约束选择暂不发布。科学结论和工程决策并不需要同一个词：证据不足可以对应保守保持旧版本，也可以对应清楚限制的小范围试用，前提是任务允许。

发布后需要确保使用的持久物与评价时相同。索引异步更新、缓存未刷新、技能包入口变化或模型路由混合，都可能让服务配置偏离候选。可以在任务记录中保存实际载入的版本，而不是只在发布页面显示一个版本号。出现故障时，这些记录帮助定位是已评价设计本身失败，还是部署条件不一致。

恢复记录也应与发布记录成对。回退到哪个版本、哪些状态恢复、哪些副作用未撤销、哪些任务已经重新评价，都应明确。恢复后的一次成功检查不能代替所有能力保持，但没有功能检查又无法确认回退实际可用。检查范围要由改变的对象和风险决定，不必给每次小改动都套同样庞大的清单。

现有工件把提示、动作、持久物与费用连接起来，读者可以沿真实事件复核协议失败、来源差异和重复计费。下一项组合比较应在这个可检查起点上增加明确干预，再用相应任务与资源记录判断它带来的变化。

\section{读取者比较与记忆输入路径}
\label{sec:ch15-reader}

\begin{sloppypar}
本节汇总第5章同一批拟合目标比较，用于说明系统级迁移读数的边界。
判决规则在第 5 章就写死了：\textbf{判决前必须同时报告换模型后的静态基线}，
否则“经验失效”与“模型更弱”不可区分。拆分 B 的完整读数如下。
\end{sloppypar}

% figures/ch15-reader.tex —— 第 15 章核心图（全书引用入口）
% 用法：在 book/chapters/ch15.tex 中 % figures/ch15-reader.tex —— 第 15 章核心图（全书引用入口）
% 用法：在 book/chapters/ch15.tex 中 % figures/ch15-reader.tex —— 第 15 章核心图（全书引用入口）
% 用法：在 book/chapters/ch15.tex 中 \input{figures/ch15-reader}
% 图主体在 ch15-reader-body.tex，便于单独编译与复用。
\begin{figure}[htbp]
\centering
\input{figures/ch15-reader-body}
\caption{E03c拆分B的两种拟合写法与两个读取者。每格评价\measured{16}题、每题\measured{5}次；格内给出回合通过率和相对该读取者静态基线的差。同读基线为\measured{0.6875}，换读基线为\measured{0.2625}。四个条件用于构造写法与读取者的交互差；区间及来源、检索与跨日期限制见正文。}
\label{fig:ch15-reader}
\end{figure}

% 图主体在 ch15-reader-body.tex，便于单独编译与复用。
\begin{figure}[htbp]
\centering
\begin{tikzpicture}[
 cell/.style={draw,rounded corners=2pt,minimum height=18mm,text width=50mm,align=center,font=\small},
 head/.style={font=\small\bfseries,align=center},
 arrow/.style={draw,dashed,gray!60,thick}
]
\node[head] at (0,2) {同读取者};\node[head] at (6.5,2) {换读取者};
\node[font=\small,rotate=90] at (-3.3,0.7) {拟合到求解器};
\node[font=\small,rotate=90] at (-3.3,-1.7) {拟合到问题域};
\node[cell,fill=blue!8] (a) at (0,0.7) {通过率 \measured{0.7000}\\相对同读基线 \measured{+0.0125}};
\node[cell,fill=orange!10] (b) at (6.5,0.7) {通过率 \measured{0.3625}\\相对换读基线 \measured{+0.1000}};
\node[cell,fill=gray!8] (c) at (0,-1.7) {通过率 \measured{0.7000}\\相对同读基线 \measured{+0.0125}};
\node[cell,fill=green!8] (d) at (6.5,-1.7) {通过率 \measured{0.4250}\\相对换读基线 \measured{+0.1625}};
\draw[arrow] (a.east)--(b.west);\draw[arrow] (c.east)--(d.west);
\end{tikzpicture}

\caption{E03c拆分B的两种拟合写法与两个读取者。每格评价\measured{16}题、每题\measured{5}次；格内给出回合通过率和相对该读取者静态基线的差。同读基线为\measured{0.6875}，换读基线为\measured{0.2625}。四个条件用于构造写法与读取者的交互差；区间及来源、检索与跨日期限制见正文。}
\label{fig:ch15-reader}
\end{figure}

% 图主体在 ch15-reader-body.tex，便于单独编译与复用。
\begin{figure}[htbp]
\centering
\begin{tikzpicture}[
 cell/.style={draw,rounded corners=2pt,minimum height=18mm,text width=50mm,align=center,font=\small},
 head/.style={font=\small\bfseries,align=center},
 arrow/.style={draw,dashed,gray!60,thick}
]
\node[head] at (0,2) {同读取者};\node[head] at (6.5,2) {换读取者};
\node[font=\small,rotate=90] at (-3.3,0.7) {拟合到求解器};
\node[font=\small,rotate=90] at (-3.3,-1.7) {拟合到问题域};
\node[cell,fill=blue!8] (a) at (0,0.7) {通过率 \measured{0.7000}\\相对同读基线 \measured{+0.0125}};
\node[cell,fill=orange!10] (b) at (6.5,0.7) {通过率 \measured{0.3625}\\相对换读基线 \measured{+0.1000}};
\node[cell,fill=gray!8] (c) at (0,-1.7) {通过率 \measured{0.7000}\\相对同读基线 \measured{+0.0125}};
\node[cell,fill=green!8] (d) at (6.5,-1.7) {通过率 \measured{0.4250}\\相对换读基线 \measured{+0.1625}};
\draw[arrow] (a.east)--(b.west);\draw[arrow] (c.east)--(d.west);
\end{tikzpicture}

\caption{E03c拆分B的两种拟合写法与两个读取者。每格评价\measured{16}题、每题\measured{5}次；格内给出回合通过率和相对该读取者静态基线的差。同读基线为\measured{0.6875}，换读基线为\measured{0.2625}。四个条件用于构造写法与读取者的交互差；区间及来源、检索与跨日期限制见正文。}
\label{fig:ch15-reader}
\end{figure}


\begin{table}[htbp]
\centering
\caption{E03c 拆分 B（标准划分）：同读取者与换读取者的两套读数。
拆分B六臂已有归档；\texttt{fit\_domain\_cross} 于 2026-09-28 续跑补齐}
\label{tab:ch15-reader}
\begin{footnotesize}
\begin{tabular}{>{\raggedright\arraybackslash}p{2.8cm}cc>{\raggedright\arraybackslash}p{4.6cm}}
\toprule
臂 & 回合通过率 & 通过 / 总运行 & 相对各自基线的收益 \\
\midrule
\texttt{base\_same}（同读基线） & \measured{0.6875} & \measured{55}/\measured{80} & （基线） \\
\texttt{base\_cross}（换读基线） & \measured{0.2625} & \measured{21}/\measured{80} & （基线） \\
\midrule
\texttt{fit\_solver\_same} & \measured{0.7000} & \measured{56}/\measured{80} & \measured{+0.0125}（\measured{$[-0.1125, +0.1250]$}） \\
\texttt{fit\_solver\_cross} & \measured{0.3625} & \measured{29}/\measured{80} & \measured{+0.1000}（\measured{$[-0.0125, +0.2125]$}） \\
\texttt{fit\_domain\_same} & \measured{0.7000} & \measured{56}/\measured{80} & \measured{+0.0125}（\measured{$[-0.1500, +0.1625]$}） \\
\texttt{fit\_domain\_cross} & \measured{0.4250} & \measured{34}/\measured{80} & \measured{+0.1625}（\measured{$[+0.0000, +0.3250]$}） \\
\bottomrule
\end{tabular}
\end{footnotesize}
\end{table}

\begin{sloppypar}
两条基线之间差 \calcnum{42.5} 个百分点；
\texttt{fit\_solver} 两侧差 \calcnum{33.8} 个百分点（\measured{0.7000} vs \measured{0.3625}），
\texttt{fit\_domain} 两侧差 \calcnum{27.5} 个百分点（\measured{0.7000} vs \measured{0.4250}）。
这些数说明静态基线与经验方案在本批任务上都有明显点值变化；它们不是对“模型”与“经验”贡献的因果份额分解。换读取者一侧两种经验方案相差\measured{0.0625}，即\calcnum{6.25}个百分点；跨日期、模型别名与任务范围的限制仍然存在。\end{sloppypar}

\begin{sloppypar}
同一模型下，P3b 比较领域拟合与求解器拟合。两者差值的点估计为 \measured{$0.0000$}，
区间 \measured{$[-0.1375, +0.1500]$}，含 0。
P3b 的跨模型读数（2026-09-28 补齐）：点估计 \measured{$+0.0625$}、
区间 \measured{$[-0.0625, +0.1750]$}，同样含 0——
点估计超过原设计参考的\calcnum{0.03}，不代表不确定性已经支持相应幅度。这里只有\measured{16}道配对题，且一次换模型不能代表跨所有读取者或跨harness。该区间未确定差异方向，也不表示两者等价。

跨模型领域拟合臂相对静态基线的区间下界显示\measured{$+0.0000$}，不能根据这个格式宣称严格为正。该臂相对静态基线的增量，与领域拟合减求解器拟合的比较，是两个不同问题；应分别使用对应差值与不确定性。\end{sloppypar}

\begin{sloppypar}
换读取者还有一笔必须记的账：它不只是准确率变了，成本也变了。
\texttt{base\_cross} 的回合计时和是 \texttt{base\_same} 的 \calcnum{2.1} 倍；
\texttt{fit\_solver\_cross} 的生成 token 是同模型那一臂的 \calcnum{5.3} 倍、
回合计时和是 \calcnum{5.6} 倍。
本次替代读取者在这些任务上通过率更低，所列方案的生成用量和回合计时和更高；这不代表所有替代模型都有相同变化。
任何只看通过率的迁移结论都漏掉了这一半。
\end{sloppypar}

若要问两种拟合写法的差是否随读取者改变，还需直接构造交互差。令 $Y_{D,C}$ 与 $Y_{S,C}$ 分别表示替代读取者下的领域拟合、求解器拟合得分，$O$ 表示源读取者，则
\begin{equation}
 I_R=(Y_{D,C}-Y_{S,C})-(Y_{D,O}-Y_{S,O}).
 \label{eq:ch15-reader-interaction}
\end{equation}
现有点值给出 $\calcnum{0.0625}-\calcnum{0.0000}=\calcnum{0.0625}$。同一读取者内，两种写法若减去共同静态基线，基线项会在相减时消去；保留目标基线仍有助于解释各方案的经验增量。计算出差的差，并没有同时得到它的不确定性。应在共同题目单位上联合构造四臂差值，保留各臂之间的配对；分别查看两侧区间、比较点值或相减两个区间端点，不能代替该交互量的估计。跨日期、模型别名和检索路径的限制也继续适用。

\subsection{从记忆文本走到真正的模型输入}
\label{sec:ch15-inputpath}

E03的离线重建提供了另一种端到端走查。审计读取归档记忆与任务陈述，重建检索结果，再到原始系统提示中检查所选文本是否确实存在。它核对的是材料进入输入的路径，不重新运行模型，也不修改历史答案。这个方法使“记忆文件里有什么”与“本题实际看到了什么”能够分开。

在拆分B中，两种拟合写法对\measured{16}道评价题的检索结果，有\measured{12}题选择的来源任务集合不同。例如对\texttt{classstate-07}，求解器拟合条件选择的是错误合同与栈状态的来源，问题域拟合条件选择的是另一个类状态来源与可变默认参数来源。归档提示核对确认相应文本确实出现，但不同条件收到的内容已不只是同一条经验换一种措辞。

这会改变 P3b 的解释。最终差值比较的是两种材料加工与检索组合；它可以描述完整方案，却没有纯粹隔离“针对读取者说话”与“描述问题域事实”这一个因素。若要更接近表达目标的作用，可以固定来源条目集合后比较写法；若要评价完整检索系统，则可以允许集合变化，但要明确将检索影响纳入效果。

同一个来源标识在不同拆分中也不一定代表同一条执行轨迹。归档的栈任务在拆分A中对应一次失败来源，在拆分B中对应一次成功来源。任务名字相同，运行、反馈和材料形成过程仍需各自绑定。若按任务名把两条来源合并，就会造成看似自相矛盾的标签；真正需要的主键应包含运行与来源事件。

原始轨迹、摘要与结构化形式还改变长度。拆分A的栈来源原始条目被截到\measured{600}字符，而摘要和结构化条目分别为\measured{51}与\measured{87}字符。较短不自动更好，较长也不自动保留完整信息；原始条目的代码可能在截断处结束，摘要则可能省去异常条件。字数、内容和检索一起变化，使简单的形式排名难以归因。

提示中出现了条目，也还差应用证据。模型可能忽略它、误解它，或者原本就会解决任务。当前记录没有可靠的逐条因果采用标记；模型自述“使用了这条经验”也不能单独证明采用。可以研究带来源标识的决策记录、受控替换或删除对照，但这些是后续设计，不能从已有提示文本中补出没有观测的内部过程。

\section{E05v2：沿真实序列审计构建与复用}
\label{sec:ch15-e05case}

本节选用已经保存的 E05 v2 工程小试作为贯穿案例。它只涉及边界错误任务族、两个训练题和复用格点零与一，没有封存评价，也没有跑完整技能系统。选择它是因为原始事件能够连成一条可核对的路径，不是因为它证明了技能收益。第8章解释机制，本节把配置、轨迹、持久物和成本接起来看。

运行标识为\texttt{E05v2-20260929T143245Z-d7a8db}。配置保存了运行代码、协议和任务的哈希；其中的\texttt{git\_head}未完整包含当时新增的执行器，仅检出该提交不足以重建运行。逐文件哈希还需与实际保留的源码分别核对。对应来源任务是\texttt{offbyone-01}，后续题为\texttt{offbyone-02}。前者修复含端点的求和，后者修复闭区间内元素计数。它们共享边界概念，但接口和输出语义不同，不能直接把一个函数当成另一个任务的答案。

在静态技能臂、一次复用格点中，来源任务通过后，系统从实际候选源文件提取代码，形成带入口签名、适用说明、来源验收题和失败条件的参考。持久物中的实现为含右端点的求和函数。准入记录为通过，来源代码进入后续提示；这个发布动作发生在来源验收之后。该验收来自旧执行路径，不能据此声称已具备操作系统隔离或不可篡改的最终评价。

参考的限制也很明确：它只在来源题上验证过，不同输入合同或接口尚未验证。下一题要求统计列表中处于闭区间的元素，系统提示把前一函数作为参考，并要求对当前题重新适配和检查。这里所谓技能复用是\textbf{提示中提供代码参考}，不是执行器直接调用一个已封装工具。把两种实现混写，会夸大这次小试已经验证的能力。

后续轨迹显示，模型依次读取当前源文件和测试文件，将严格不等号改为含端点的不等号，写入成功后调用测试，再结束回合。最终检查通过。序列记录的 \texttt{retrieved} 为真，实际系统提示中可见代码参考，当前题也通过验收。该字段在这条路径上标记参考注入，没有单独执行候选查询或排序。

但这两项记录还不能证明参考造成了通过。当前测试可读，模型也可能根据题意与测试独立发现边界错误；无参考的重规划臂在本次相应题目上同样通过。轨迹中没有一个能够隔离参考因果作用的反事实。完整走查的价值是确认机制通路，而不是把通路存在当成收益已经成立。

来源验收也没有穷尽代码语义。通过当前有限测试，不能保证求和函数处理任意类型或异常合同都正确。契约字段写明限制是良好的记录，但字段存在不代表约束已经被自动执行。第8章所要求的适用性检查和回归仍需要具体实现与评价，不能从一份格式完整的参考对象中推导出来。

\subsection{同一条路径上的调用与令牌}

\begin{sloppypar}
这条静态序列含\measured{2}个主回合，构建使用\measured{6}次模型调用，后续题使用\measured{6}次，总计\measured{12}次。提示令牌为\measured{5365}，生成令牌为\measured{358}，两者合计\calcnum{5723}。回合计时和约\measured{11.386}秒，但它不是整次试跑的完整墙钟，更不是货币费用。
\end{sloppypar}

\begin{sloppypar}
同格点重规划臂也是\measured{12}次调用，提示和生成令牌分别为\measured{4753}与\measured{375}，合计\calcnum{5128}。静态技能臂在这条序列上多\calcnum{595}个令牌；调用数相同，并未显示调用节省。只有一次短序列，不能据此拟合长期斜率或给出稳定回本点。
\end{sloppypar}

比较时不能把构建费在每道后续题上再加一次，也不能将已经累计的复用调用乘以复用次数。这里的序列总账已经包含构建和后续实际调用，角色拆分用于解释总量，不是另收一遍费用。每个格点又独立重建状态，较小格点的回合不是较大格点免费继承的前缀。

\subsection{失败的构建为什么仍然是重要结果}
\label{sec:ch15-rejectedbuild}

动态技能臂在同一格点走出了另一条路径。来源回合先列出文件，随后读取源代码与测试；模型写出了修复方向，却在工具调用文本中使用了不符合当时单行协议的结构。记录中出现“没有可解析的调用或结束标记”的错误，之后来源验收未通过。这里最先暴露的是协议兼容性，而不是算法已经证明边界修复不可能。

有两个细节不能忽略。其一，某一轮回复包含多个调用意图，但旧解析器只执行匹配到的一个调用；不能按模型文字中出现几次工具名，就计几次实际工具执行。其二，失败的写入意图没有对应成功写入事件，因此不能把意图中的代码当成最终文件版本。成本、状态与最终判定必须依据实际执行事件，而不是依据模型叙述。

该来源回合花了\measured{4}次模型调用，序列将准入原因记为来源失败，技能没有发布。下一道计数题仍继续执行并通过，使用\measured{6}次调用；但它的检索标记为否，序列最终技能为空。因此可以说这条后续任务通过了，不能说动态技能迁移成功了，因为该序列根本没有可供迁移的技能。

这也是为什么本章不把臂名称当成事件事实。一个名为动态技能的条件可以经历空库、构建失败和普通执行；只有记录显示候选被准入、被取出并进入实际使用路径时，才能讨论对应机制是否发生。若只汇总臂名称和成功率，会把不同状态混在同一解释里。

动态序列合计\measured{10}次调用，比相应静态与重规划序列少\calcnum{2}次。少花调用来自构建提前失败，不能命名为效率提升。把费用最少的条件直接选成胜者，会奖励没有完成目标的系统；因此费用比较必须同时呈现完成质量和中止原因。

本次没有发生修订回合，修订次数为\measured{0}。它表示这次记录中没有走到相应修订路径，不表示修订器成功率为零，也不表示修订预算机制已经在真实失败复用中得到充分检验。动态臂的名称不提供分母，修订效果需要实际修订尝试。

另一零复用格点的动态来源回合也因格式问题失败，使用\measured{2}次调用。整个工程小试共有\measured{9}个主回合、\measured{7}次通过、\measured{48}次模型调用，汇总记录的本进程后端重试为\measured{0}；这不是独立的服务端计费核验。所有这些数描述一次有限工程运行；它们不能形成 P4 所要求的质量保持、充分复用与成本判据。

\subsubsection{一次验收动作与最终判定也要分开}

来源失败轨迹中的\texttt{test\_runs}为零，表示模型没有成功调用回合内的测试工具；轨迹末尾仍有运行器执行的最终验收判定。两者属于不同记录位置，不能把工具计数为零写成没有进行任何验收。反过来，回合内一次测试通过，也不应代替结束时对实际最终文件的核对，因为之后可能继续发生修改。

因此复核时应同时看模型动作、工具结果与终局判定来源。模型提出写文件请求，工具确认写入成功，最终验证器检查当前文件，是三个连续但不同的事件。缺少中间确认时，不能仅凭请求中的正确代码宣布环境已经改变；缺少最终判定时，也不能把工具没有报错当成任务通过。

这条细节解释了为什么原始轨迹值得保存：汇总表只留下一个失败标记，看不见失败发生在动作协议、环境执行还是任务语义。逐事件记录允许后续修复对应环节，同时保留旧运行实际完成范围。修复后再验证，是新的观察，不会追溯改变这次来源失败的身份。

\subsection{同一份案例中，三种分母各回答什么}
\label{sec:ch15-denominators}

E05工程小试的九个主回合并不对应九道独立任务。三个条件分别在零复用格点运行一题、在一次复用格点运行两题，因此有 $\calcnum{3}\times(\calcnum{1}+\calcnum{2})=\calcnum{9}$ 个主回合；独特题目只有\measured{2}道，任务家族只有\measured{1}个。回合数、题目数和家族数都是真实计数，但对外推范围的含义不同。

按主回合汇总，七次通过可以描述这次装置运行完成了哪些检查；按题目比较，可以查看同一个求和或计数问题在不同条件下发生了什么；按家族推断，则只能承认没有多个家族支持广泛迁移。把九个回合当作九种不同问题，会让样本覆盖显得比实际更大。

两个格点也不是一条越来越长的连续运行。每个格点重建独立状态，来源题会重新执行。零格点的一次来源成功，不会自动成为一次复用格点的免费技能来源。这样可以避免不透明地共享状态，但相应构建费用会真实发生多次，必须按项目总账保留。讨论某一部署序列的费用时，则只取该序列内的构建与使用。

第三种分母是机制机会。一次复用格点提供了一个后续任务位置，却不保证存在已准入技能。动态臂来源失败后，该位置仍发生任务执行，但没有注入代码参考。可以同时报告计划复用机会和实际注入次数，分别回答装置安排了多少机会、参考实际进入提示多少次。不能用计划机会代替实际注入，也不能把注入次数当作采用次数。

在本例中计算“注入参考后通过率”，应限定到实际注入参考的回合，但这会选择一个受来源成功影响的子集。它适合诊断已进入机制的路径，却不能替代完整方案的总体表现，因为构建失败被排除了。总体方案比较与条件子集诊断可以并列，不能挑更好的一项作为唯一结果。

修订也有同样的问题。修订机会需要先发生相应失败，再检查预算和准入条件。本次修订尝试数为零，因而没有修订成功率的可用分母。把零除以零填成零或百分之百都不正确。报告“未发生修订尝试”可以准确表达事实，也提醒下一项测试应该覆盖哪条路径。

\subsection{把调用差、质量差与回本点放在同一页}
\label{sec:ch15-costqualitypage}

本次静态技能与重规划在两个格点的调用数都相等：零格点各\measured{6}次，一次复用格点各\measured{12}次。程序可以据此报告观察格点上的调用差为零，但“平价”不等于“节省”，更不等于已经找到持续复用后的长期回本规律。费用曲线需要更多使用机会，质量约束还要单独验证。

令 $\Delta C(k)$ 表示技能方案减去相同格点重规划方案的总费用。观察到某个格点 $\Delta C(k)\leq0$，只能说明这项费用口径在这个格点不更高。如果提示令牌的差仍为正，调用平价也不代表所有资源平价；如果后续格点再次变正，第一次交点也不是永久收益边界。第8章的成本模型因此必须与实际格点和质量范围绑定。

对于动态序列，调用差虽然为负，却伴随构建失败，无法作为满足质量条件的省费点。可以将这类记录画在成本与质量的二维图上，显示它为什么不满足预定接受要求，而不是把失败记录删掉。删除会掩盖机制故障，只按成本排名又会奖励故障；同时呈现两项才保留完整解释。

质量差也要选择对应指标。若主张是后续任务首轮完成，来源题通过率不是同一指标；若主张是整个工作流程完成，来源构建失败就必须进入结果。相同数据可以支持不同定义的描述，但不同定义不能在方法之间切换。否则一条路线按最好看的后续表现计分，另一条按完整流程计分，比较会失去共同基础。

现有记录覆盖短序列的构建与一次复用，长期复用、漂移后修订、多家族外推和货币费用仍为空项。下一份比较应从其中选择会改变决策的问题，补上相应的任务与资源记录。

当后续新版本运行到来时，应新增一页对应结果，保留旧页作为工程历史。版本变化可能改善协议稳定性，也可能改变成本与行为；只有实际运行才知道。将旧失败留在正确身份下，可以让改进过程本身成为可复核案例，而不是只留下最终一次最顺利的故事。


% ---------- 9. 条件变化 ----------
\subsection{修复实现之后，为什么不能改写旧轨迹}
\label{sec:ch15-versionseparation}

发现格式兼容问题后，可以修复协议处理并做离线检查。但修复后的执行器面对相同提示，可能执行不同动作，后续观察与模型回复也可能变化。旧轨迹只是旧实现实际走过的路径，不能把其中失败判定替换为新解析器“本来会成功”的判定，再称为新版真实运行。

可以离线重放文本以研究解析器行为，但其结果身份应写成程序检查：给定同一段历史回复，新旧解析器分别输出什么。它没有重新取得模型回复，也没有执行整个下游环境，更不能估计新版本整体通过率。若要评价新设计，需要独立运行标识与对应的实际轨迹。

本书因此保留 v2 的原始失败记录，将 v3 的改动和离线验证另行冻结。v3尚未获得本轮真实结果，不能用修复完成来替它填写效果。历史运行版本、代码哈希和设计身份一起保存，便于将来比较，但不会把不同版本的半截序列合并。

费用也遵守同一边界。旧版已经发出的请求是实际支出，修复不能让它们消失；新版若再运行，会产生另一笔费用。某个请求失败且服务端未返回用量时，应保留未知费用，而不是按未成功取得结果就记为零。本次已保存的 v2 审计记录了已知调用与令牌，美元费用仍未知。

这一纪律也适用于跨日期续跑。E03c补齐时必须保留原配置与续跑事实；最终十臂齐全解决的是缺失臂问题，不会消除模型服务随时间变化的可能。将不同时间的数据视为同一研究，可以有明确协议，但不能为了表格整齐而抹掉时间标识或别名限制。

\subsection{把案例包组织成读者能复核的顺序}
\label{sec:ch15-casebundle}

一个可读案例包可以按五个入口组织。先看设计与任务清单，确定被比较的条件；再看逐回合轨迹，确认实际动作；随后看序列事件，确认准入和复用；接着看唯一费用账；最后才看汇总和解释。这个顺序让读者可以发现哪一层出现缺失，而不是先接受一个结论再寻找支持。

本节案例的原始配置、六条序列完成记录及九条轨迹均位于上述运行目录。新增离线案例索引只记录路径、哈希和从原字段重算的数字，不改原文件。哈希用于确认文件身份；它不能证明任务足够独立、检查器语义正确或分析无偏。工件完整性和研究有效性是两项不同验收。

汇总中的每个数字都应能够回到分母。例如通过次数来自哪些主回合，修订是否另算，构建失败是否仍保留，重试如何计数。若正文写的是后续复用效果，就不能悄悄把来源构建成功也放入分子；若写的是整个序列成本，则不能只取后续阶段。选择分母是研究问题的一部分。

案例包还应收录未发生的关键路径。本次没有修订、没有真实恢复、没有封存评价，这些是覆盖缺口；不能把它们记成零失败或百分之百安全。可以在下一项设计中针对这些路径增加检查，但结果到来之前，缺口依然存在。将缺口写在案例旁边，比在全书末尾统一加一句“尚有局限”更方便复核。


\section{从既有迁移范围设计新的检验}
\label{sec:ch15-family}

\begin{sloppypar}
跨族评价改变了任务覆盖范围。E03拆分A将两个任务族全部留出。三个记忆臂保存的判定计数为原始轨迹\measured{71/80}、摘要\measured{62/80}和结构化\measured{61/80}；这允许描述所测完整方案，却不能归因为加工深度。
\end{sloppypar}

\begin{sloppypar}
第~\ref{ch05}章\S\ref{sec:ch05-record-conflict}集中说明了该拆分的基线记录差异及采用边界。本节保留三个记忆臂的观察和不涉及该基线的直接比较。拆分B的任务子集不同，其基线不能移作拆分A的参照。
\end{sloppypar}

\begin{sloppypar}
结构化减摘要的点估计为\measured{$-0.0125$}，历史题级区间\measured{$[-0.1500,+0.1125]$}跨零，方向未定，不等于二者等效。P3进一步问形式与容量怎样共同影响迁移；这需要交叉操纵两个因素，并按目标范围安排独立留出族，不能从现有形式比较直接推断。\end{sloppypar}

\begin{sloppypar}
跨代码仓库这一步本书没有做。任务集 v1 的 \measured{64} 道题全部是同一类形态：
短小 Python 源码与单元测试验收器。换一个仓库可能改变验收器形态、依赖或失败模式，
那是外部效度问题，见下一节。
\end{sloppypar}

% ---------- 6. 跨域迁移：本书为什么没做 ----------
\subsection{从代码任务到环境操作：设计新的检验}
\label{sec:ch15-domain}

\begin{sloppypar}
本书的全部实测来自一个环境：可执行的编程任务，验收器是单元测试。
这个选择的好处是验收便宜、可重复、可封存；
代价是它主要覆盖由代码断言检查的结果，包括返回值与局部状态，不能代表所有真实环境目标。
\end{sloppypar}

\begin{sloppypar}
另一类任务长得很不一样：工具操作类任务要在环境里留下一串副作用，
判断对错靠\textbf{环境的最终状态}，而不是一次返回值。
提纲里登记的迁移环境是这一类（AppWorld，2407.18901），
本书首版\textbf{没有运行它}；现已核验原论文的状态评价机制\cite{appworld-2407}——
本节说明扩展到这类环境需要检验什么，不给任何数字。
\end{sloppypar}

\begin{table}[htbp]
\centering
\caption{内部效度与外部效度：本书做到了哪一步}
\label{tab:ch15-validity}
\begin{footnotesize}
\begin{tabular}{>{\raggedright\arraybackslash}p{2.6cm}>{\raggedright\arraybackslash}p{4.4cm}>{\raggedright\arraybackslash}p{4.4cm}}
\toprule
效度 & 要回答的问题 & 本书的状态 \\
\midrule
内部效度（同域） & 收益确实来自被改动的那个组件吗 & 已有四格、族留出与换读取者记录；仍有预算与来源混淆 \\
\midrule
外部效度（跨域） & 换一类答案形态后还在吗 & \textbf{未做}：迁移环境未运行，无数字 \\
\midrule
外部效度（公开基准） & 在别人造的题上还在吗 & \textbf{未做}；且存在污染风险（见下） \\
\bottomrule
\end{tabular}
\end{footnotesize}
\end{table}

\begin{sloppypar}
公开基准这一侧还有一个不能绕开的风险：\textbf{无法排除基础模型在预训练时已经见过公开数据}。
自建新题减少直接复用既有题面的一部分风险，却不能排除常见模板、解法或生成来源与模型知识重叠。公开基准也不能一概被称为分数虚假；应分别审查污染、覆盖与评价边界。本书自建集的主要已知限制是范围窄，不能据此声称已经排除全部训练数据重叠。
\end{sloppypar}

\begin{pitfallbox}
本书的负向点估计不能直接替公开基准作结论。
“结构化记忆减摘要的点估计为负、区间跨零”这条记录的适用范围是：
任务集 v1、\measured{16} 道封存题、\measured{5} 次运行、\texttt{deepseek-flash}、单元测试验收器。
改变这些条件后是否仍有同样读数，需要新证据。
\end{pitfallbox}

% ---------- 7. 三类更新案例 ----------
\subsection{迁移是一组轴，不是一次换环境}
\label{sec:ch15-transferaxes}

本书已经包含同族新题、有限族留出和替代读取者的观察。它们改变的是不同轴：题目实例、失败模式范围和执行模型。跨仓库会进一步改变代码结构与依赖；换执行框架会改变工具协议和上下文组织；转向操作型环境则改变状态与副作用。将这些统称为迁移，便于概括，却不足以设计或解释比较。

例如，一段记忆写着“先运行这个测试命令，再检查这个文件”。它在同仓库新题中可能适用，跨仓库时文件与命令不存在；若重写为更一般的失败诊断规则，表示更抽象，但也可能失去执行细节。哪种形式更有用取决于目标范围，不能预先规定抽象越高迁移越好。E03 的形式比较只覆盖自己的处理和任务条件。

换模型也不只是换阅读理解能力。相同文本可能被不同分词器计成不同长度，工具调用格式稳定性、上下文利用和结束习惯也可能不同。如果后端适配层跟着改变，观察到的是模型与适配条件的联合变化。可以尽量统一接口并保存实际输入输出，但不能仅凭请求字段不同就知道具体机制。

换框架则可能在基础模型不变时改变结果。检索到的经验放在系统提示还是工具返回中，失败后是否自动重试，工具异常是否转成文字反馈，都会改变策略所见观察。一个经验条目在框架甲有效，在框架乙失效，未必是条目内容错误；也可能是没有送到执行者、被截断或失去了适用标识。迁移检查需要追踪材料进入实际调用的路径。

外部环境的评价应按轴逐步扩大。可以先固定任务家族更换仓库，再固定仓库更换读取者，最后研究二者同时变化。这样做不是保证没有交互，而是降低最初解释的难度。如果实际部署必须同时换多项，也可评价整个迁移方案，只需避免把总变化归给单一轴。

每项迁移还应保留目标环境自己的静态基线。源环境基线说明出发点，目标环境基线说明不携带新增经验时系统能做到什么。两者都能帮助理解结果，但不能互相替代。第15章开头的读取者例子之所以容易误读，正是因为绝对表现与经验增量使用了不同参照。

\subsection{从返回值检查走向环境状态检查}
\label{sec:ch15-stateevaluation}

AppWorld原论文提供了一个可操作的状态评价机制：比较初始与最终数据库的差分 $\Delta D$，检查必须发生的变化是否齐全，同时排除允许范围之外的附带修改\cite{appworld-2407}。沿用其集合记号，可将这两项要求合写为
\[
 C_{\mathrm{expected}}\subseteq\Delta D\subseteq C_{\mathrm{expected}}\cup C_{\mathrm{allow}}.
\]
这并不要求所有路径复制参考操作；具体测试还可以查询最终数据库中的关联属性，不能只检查差分里出现的字符串。

该文的TGC按通过全部验收的任务聚合，SGC则按同一场景下全部任务变体都通过的场景聚合。一个变体失败时，其他变体通过也不能让该场景进入SGC成功分子；这是对变化要求与初始状态的检查，与第\ref{ch14}章同一任务随机复测的 $\mathrm{pass}^{k}$ 不同。本书没有运行该基准，这里采用的是评价机制，不把原文模型成绩计入本书迁移结果。

下面另设一个本书教学任务：在模拟日程中，把指定会议改到用户要求的时间，同时保留参与者与其他会议。目标不只是输出一句“已修改”，而是环境中的指定记录满足要求。若系统改错另一场会议，即使最后回答语气肯定，也没有完成任务；若正确改了目标却删除了其他会议，也不能只按目标字段成功判定。

可以把要求拆成必须变化、允许变化和必须保持三部分。必须变化是目标会议的时间；允许变化可能是与操作一致的更新时间戳；必须保持包括会议身份、参与者及未指定记录。具体字段由任务规范决定，不能让智能体自行决定哪些副作用可以忽略。这个分类是评价设计，不是对 AppWorld 任一具体任务的复刻。

多条操作路径都可能达到同一个有效状态。系统可以先查列表再定位会议，也可以使用明确标识直接读取；只要权限与任务要求允许，评价不必要求逐步复制参考轨迹。另一方面，仅检查最终状态可能漏掉曾经发生又被撤销的违规操作。若过程约束重要，应另有不可由候选改写的动作记录，最终状态与过程证据各自负责一部分要求。$\tau$-bench原文给出一个具体边界：即使最终退货状态和必要答复都正确，未经用户明确确认就执行退货仍可能违反策略；它因而把所用结果奖励视为成功的必要而非充分条件\cite{tau-bench-2406}。

环境初始状态必须恢复到同一条件，尤其当多个方法处理同一任务时。如果第二个方法开始时目标已经被第一个改好，就可能无需正确行动也得到成功。恢复不仅是数据库文件，还可能包括时间、交互状态和缓存。可重建的模拟环境有利于比较，但恢复是否完整仍应通过实际检查确认。

结果反馈的可见范围也会改变学习问题。若每一步都告诉智能体哪个验收条件失败，它获得了丰富的开发信息；若只在结束后判定成功，它面对的是稀疏反馈。两种条件都可以研究，但不能将前一种的更新效果推广为没有反馈时也能达到。测试时使用了最终评价信息，就应明确报告，而不是继续称评价完全隐藏。

\subsection{在操作型环境里，什么才算撤销}
\label{sec:ch15-sideeffects}

第~\ref{ch14}章区分工件恢复、行为恢复和外部补偿。操作型环境还需逐字段核对：日程案例应恢复会议时间、参与者和其他受影响记录；随后确认错误经验已被撤下，执行器加载了预期版本。模拟环境可以复位，并不意味着真实操作从未发生。

迁移评价因此至少保存初始状态、实际动作、最终状态、补偿记录和恢复后的读取检查。E12新版的受限记忆设计不覆盖这些操作型环境条件，也尚无本书真实恢复率；这里保留的是下一项设计的检查对象。

\section{用审计结果选择下一项干预}
\label{sec:ch15-errorrouting}

完整案例的错误谱最好与处理动作相连。若模型没有找到相关记忆，应先检查索引和查询；若找到了却把适用条件读反，应检查表达与上下文；若推理正确但工具调用不能解析，应检查协议；若工具执行成功但目标未完成，应检查任务语义与评价覆盖。它们都表现为最终失败，却不需要同一种修复。

E05动态构建就是协议路径的例子。模型文本已经指出右端点问题，但写入动作没有执行。继续给它补充更多边界知识，未必能解决单行调用协议与实际输出不兼容的问题。此时更有信息量的工作是离线复现解析行为、完善共同协议处理，再用新的受控运行检验变化；不是把失败直接归类为技能能力不足。

E03检索来源集合不同则属于材料路径。若一条路线没有收到目标相关条目，不能只通过润色生成器来修复；需要检查检索分数、条目选择与注入限制。相反，固定检索集合后仍出现稳定差异，才更有机会研究表达与读取方式。错误谱帮助决定下一次干预放在哪一层，而不是替代那次干预的证据。

对多组件系统，可以保留主要故障分类和不确定原因。一个回合可能同时有错误检索与工具异常，不必强行选出唯一根因；但也不应把所有可能性都勾选。记录直接观测到的失败、支持该判断的事件，以及需要额外检查的假设，能让后续修复更有针对性。

故障分类规则应有版本。若新发现一种解析失败类型，旧数据可以在保留原始记录的前提下重新标注，但应说明这是分析修订。分类改进可能改变各类比例，不代表系统本身的行为发生了变化。把分析版本与运行版本分别保存，可以避免用新版分类结果制造一条假的能力改善曲线。

还需防止只挑容易解释的失败。可先按预定抽样或完整范围检查，再展示有代表性的案例；若案例是在看完结果后挑选，应如实标为说明性案例。本章的E05序列属于后一种用途：它展示可复核机制，不估计这种失败在其他任务族中的发生率。

\subsection{把已有案例转换成下一份比较方案}
\label{sec:ch15-nextcomparison}

若下一次资源只够解决一个问题，应先选择会改变解释的缺口。对E05，可以先确认新协议是否让来源构建与复用路径按设计发生；这仍是工程验证。对E03的写法比较，可以固定来源集合与材料预算，降低检索变化带来的混淆。对跨域主张，则需要新的环境与评价范围，不能仅增加现有短程序的重复次数。

新方案应在运行前写清主结果。若关心技能复用后的首轮表现，来源构建成功和修订后成功应分别报告；若关心完成任务的总资源，则构建、失败与修订都进入费用。一个序列可以支持多个描述，但主判据需预先明确，不能根据哪张表好看才决定主要结论。

样本规划也要跟随目标单位。增加同一任务的重复，可以更好观察调用随机性；增加同一模板变体，可以扩大实例覆盖；增加独立任务家族，才更接近跨机制范围。三者不可用一个总回合数替换。任务家族数量有限时，可以收窄主张或报告探索性结果，但不能把相关回合当成完全独立样本来制造精确区间。

如果后端或模型身份在运行中改变，应先保存已完成部分和时间边界，判断原对照是否仍然成立。必要时将其标为不同批次，分析时解释变化。保持所有数据不丢失，与强行把所有数据合成同一个主结果，是不同要求。前者有助于复核，后者需要设计依据。

对于预算已耗尽但结果未定的试验，正确的交付是实际完成范围、已用资源、缺失部分和下一步设计。没有必要为了让案例看起来闭环而补写一个成功或失败判决。第16章将这些设计、覆盖和精度问题组织为研究议程，区分已有观察与下一项对照。


\section{条件变化：这套案例什么时候不适用}
\label{sec:ch15-conditions}

\begin{itemize}
  \item \textbf{持久物不可逐条解释时}：仍可在接口兼容的读取者间做受控比较，但需要记录表示与解码接口，不能要求它具有自然语言条目的相同审计字段。
  \item \textbf{读取者比写入者强时}：本节的方向可能反过来。
        本书只测了一个方向（换到另一个模型），没有测“换到更强的模型”。
  \item \textbf{验收器形态改变时}：从单元测试换成环境最终状态检验之后，
        “通过 / 未通过”的含义变了，可以并列描述，但不能不说明指标和任务变化就直接合并或作因果比较。
  \item \textbf{任务族之间高度相关时}：名义留出不保证机制新颖，应说明共享模板、解法与来源；共享部分技能也可能正是希望测量的迁移，不能一概宣布留出无效。
  \item \textbf{读取者与写入者共享权重时}：提示、工具与harness仍可能不同；需按实际改变条件解释比较，不能仅凭共享权重推断错误完全相同。
\end{itemize}

% ---------- 10. 能说与不能说 ----------
\label{sec:ch15-claims}

已有记录支持三种后续干预。对读取者迁移，应保留目标读取者自己的静态基线，并固定来源与信息范围；对表达目标，应在共同来源集合上改变写法；对技能成本，应先验证构建、准入与复用路径，再增加独立任务族和复用机会。一次绝对分数下降、一个正向点估计或一条成功轨迹，都不能替代这三类比较。

拆分A按第5章说明的采用边界阅读，拆分B的区间尚未确定拟合目标差异；E05的提前失败则优先指向协议兼容性检查。将它们分别送回对应的设计环节，比把全部不确定结果统一归为模型能力不足更有操作意义。

\begin{principlebox}
本章的经验收益都需要说明处理结果与对应基线的差。
换读取者或任务范围之后，差值可能变化；预算重分配也可能改变比较结果。
这些都需要实际对照，不能把“可能变化”写成每次必然反号。
所以“收益”这个词在本书里从来不是单个数字，而是\textbf{数字 + 它减去的那个基线 + 减的时候用的预算}。
三者缺一，剩下的那个数就可以被读成任何结论——
本章开头的 \measured{0.3625} 就是这么被读成两种相反结论的。
\end{principlebox}

\falsifier{P3b需要针对原拟合目标主张的有效比较。差异区间跨零只表示未定，不是证伪条件。应保留对应静态基线、来源一致性、目标读取者和harness范围，并预定效应方向及判断规则；当前同模型与跨模型比较均不作判决。}

\openquestion{本书只观察了一个替代读取者。下一项迁移比较应在冻结来源、任务和执行接口之后，明确增加哪些读取者或任务族，按相应独立单位规划精度，并同时保留目标基线与完整资源范围。跨日期续跑补齐了历史记录，但不会自动增加独立读取者数量或消除服务时段差异。}

案例索引还保留了原始文件哈希，便于区分后续解释修订与原始数据变化。解释可以更正，历史轨迹的身份应保持可追踪。

\section*{习题}
\begin{enumerate}[label=15-\arabic*.,ref=15-\arabic*]
  \item\label{ex:ch15-1} 表~\ref{tab:ch15-reader} 里 \texttt{fit\_solver\_cross} 的 \measured{0.3625}
        分别减去 \measured{0.6875} 与 \measured{0.2625} 得到不同符号。
        请说明为什么“同时报告换读取者后的静态基线”这一条纪律
        不能替换成“同时报告换读取者前的基线与经验臂”。若要比较两种写法的读取者交互，应如何构造式~\ref{eq:ch15-reader-interaction}中的四格差值？
  \item\label{ex:ch15-2} 拆分A的汇总与保存轨迹给出不同基线计数。请说明为什么不能直接选一个数纠正正文，
        并列出需要暂停的基线依赖量和可以继续复核的直接比较。
  \item\label{ex:ch15-3} 表~\ref{tab:ch15-validity} 有两行标“未做”。
        请为其中跨域那一行的迁移环境写出一份预注册清单：
        要报哪几个量、样本量按什么定、以及什么结果算“结论不变”；另请区分任务完成、场景完成与过程合规。
  \item\label{ex:ch15-4} 换读取者使生成 token 涨到 \calcnum{5.3} 倍、回合计时和 \calcnum{5.6} 倍。
        若同一方案再次遇到这样的资源增量，而总预算固定，可以有几种分配办法？
        每种付法会改变哪一个可比性条件？
  \item\label{ex:ch15-5} 相对各自静态基线的题级收益点估计，在同读时为正、换读后为负的题数为：\texttt{fit\_solver} \measured{1}/\measured{16}道，
        \texttt{fit\_domain} \measured{2}/\measured{16}道；每题在每个条件下运行\measured{5}次。
        请说明这些描述案例怎样为第5章的双字段设计提供动机，为什么还不能证明真实效应反号或该设计唯一必要，
        以及要把它升级为可判决的命题还需要什么。
\end{enumerate}

\section*{常见误解}
\begin{itemize}
  \item \textbf{“静态基线就是没有 Agent”。}
        本书的基线保留任务工具，但不发布跨任务经验；
        用不同性质的基线会改变问题，不能在没有对应实验时声称一定夸大某个固定幅度。
  \item \textbf{“拆分 A 与拆分 B 的差就是记忆效应”。}
        两套拆分的任务子集不同；拆分A基线另有记录冲突，暂停采用。
        有效基线内的差才回答对应记忆增量，不能借拆分B绕过这一缺口。
  \item \textbf{“换模型后分数掉了就是经验失效”。}
        绝对表现差与相对基线的经验增量不同；
        不报换读基线，这句话不可判。
  \item \textbf{“迁移检验只要报通过率”。}
        换读取者同时抬高成本（生成 token \calcnum{5.3} 倍、回合计时和 \calcnum{5.6} 倍）；
        只看通过率的迁移结论漏掉一半。
\end{itemize}

\section*{进一步阅读}
\begin{itemize}
  \item \cite{coding-survey-2608}：§5.2 把组件退役与回滚成功率列为应当报告的度量，
        这是该综述建议的评价方向，不能改写为所有已有工作都没有报告；本书自身尚无真实恢复率。
  \item \cite{harness-eval-2607}：核对外部任务上的rollout预算、选择方式与留出测试，
        是本章等预算四格的参照物；读它时对照第~\ref{ch03} 章四格的划法。
  \item \cite{anatomy-memory-2602}：表4同时列出答题分与记忆操作中的可恢复格式偏差。
        其中主干参与答题和记忆维护，读写条件一起变化。对照本章可检查为何迁移设计要分别固定写入与读取条件，并保留目标读取者的静态基线。
  \item 附录~\ref{app:lab}~\S\ref{sec:appb-tasks}：任务集 v1 的族划分与两套拆分的落点。
  \item 附录~\ref{app:cost}：换读取者那笔成本账的实物量口径与记账字段。
\end{itemize}
